%========================================================
% DEDICATION PAGE
%========================================================
\newpage
\thispagestyle{empty}
\vspace*{\fill}
\begin{center}
	{\large \textit{Dedications}}\\[2cm]
	
	\textit{To our families, for their unwavering support and sacrifices throughout this journey.}\\[1cm]
	
	\textit{To our teachers, who have transmitted far more than knowledge.}\\[1cm]
	
	\textit{To our supervisor, Dr. Zakaria BENZADRI, for his guidance, availability, and valuable advice.}\\[1cm]
	
	\textit{To all those who believe that technology can be a driving force for a better future.}\\[1cm]
	
\end{center}
\vspace*{\fill}

%========================================================
% ACKNOWLEDGEMENTS PAGE
%========================================================
\newpage
\thispagestyle{empty}
\chapter*{With Gratitude}
\addcontentsline{toc}{chapter}{Acknowledgements}

We would like to express our sincere gratitude to all the people who contributed, directly or indirectly, to the completion of this work.

Our deepest thanks go to our supervisor, \textbf{Dr. Zakaria BENZADRI}, for the trust he placed in us by proposing this topic, for his valuable guidance, constant availability, and the scientific rigor he instilled in us throughout this project.

We also extend our appreciation to the faculty members of the \textbf{Department of Software Technologies and Information Systems} at Constantine 2 University for the quality of education provided during our academic journey.

We would also like to thank the members of the jury for accepting to evaluate this work and for their valuable remarks that will contribute to its improvement.

Finally, we express our heartfelt gratitude to our families for their patience, encouragement, and moral support throughout our years of study.
%========================================================
% ABSTRACTS — place BEFORE \tableofcontents
%========================================================

% ---- ENGLISH ABSTRACT ----------------------------------
\newpage
\thispagestyle{empty}
\chapter*{Abstract}
\addcontentsline{toc}{chapter}{Abstract}

The Algerian agricultural sector continues to face persistent structural challenges, including over-intermediation, price opacity, and inefficient logistics coordination between producers, buyers, and transporters. This report presents \textbf{AgriGov Market}, a national web-based platform designed to digitalize and regulate direct agricultural trade under the official supervision of the Ministry of Agriculture.

The system connects four primary actors — Farmers, Buyers, Transporters, and the Ministry — within a unified digital ecosystem. Farmers can list and manage their products, coordinate logistics, and track orders; Buyers can browse a structured marketplace, place orders, and consult their purchase history; Transporters can accept and execute delivery missions via a mobile-first interface; and the Ministry retains full administrative control over user validation, product categorization, and official price regulation.

The development process followed a standard software engineering methodology structured in two phases: a Requirements Analysis phase, which produced use case diagrams, system sequence diagrams, and interface mockups; and a System Design phase, which yielded a class diagram, a global activity diagram, a relational database schema, and a deployment architecture. The platform was implemented using Django (REST API), Next.js (web frontend), React Native (mobile), and PostgreSQL hosted on Neon.

\vspace{0.5cm}
\noindent\textbf{Keywords:} agricultural platform, digital marketplace, direct trade, logistics, price regulation, Django, Next.js, React Native, Ministry of Agriculture, Algeria.

% ---- FRENCH ABSTRACT -----------------------------------
\newpage
\thispagestyle{empty}
\chapter*{Résumé}
\addcontentsline{toc}{chapter}{Résumé}

Le secteur agricole algérien continue de faire face à des défis structurels persistants, notamment la surintermédiaiton, l'opacité des prix et la coordination logistique défaillante entre producteurs, acheteurs et transporteurs. Ce rapport présente \textbf{AgriGov Market}, une plateforme nationale en ligne conçue pour numériser et réguler le commerce agricole direct sous la supervision officielle du Ministère de l'Agriculture.

Le système connecte quatre acteurs principaux — les Agriculteurs, les Acheteurs, les Transporteurs et le Ministère — au sein d'un écosystème numérique unifié. Les agriculteurs peuvent référencer et gérer leurs produits, coordonner la logistique et suivre les commandes ; les acheteurs peuvent parcourir un catalogue structuré, passer des commandes et consulter leur historique d'achats ; les transporteurs peuvent accepter et exécuter des missions de livraison via une interface orientée mobile ; et le Ministère dispose d'un contrôle administratif complet sur la validation des utilisateurs, la gestion des catégories de produits et la régulation officielle des prix.

Le développement a suivi une méthodologie d'ingénierie logicielle structurée en deux phases : une phase d'Analyse des Besoins, qui a produit des diagrammes de cas d'utilisation, des diagrammes de séquence système et des maquettes d'interface ; et une phase de Conception du Système, qui a abouti à un diagramme de classes, un diagramme d'activité global, un schéma de base de données relationnel et une architecture de déploiement. La plateforme a été implémentée avec Django (API REST), Next.js (frontend web), React Native (mobile) et PostgreSQL hébergé sur Neon.

\vspace{0.5cm}
\noindent\textbf{Mots-clés :} plateforme agricole, marché numérique, commerce direct, logistique, régulation des prix, Django, Next.js, React Native, Ministère de l'Agriculture, Algérie.

% ---- ARABIC ABSTRACT -----------------------------------
\newpage
\begin{Arabic}
	
	\chapter*{الملخص}
\raggedleft



لا يزال القطاع الفلاحي الجزائري يواجه تحديات هيكلية مستمرة، من بينها كثرة الوسطاء، غياب شفافية الأسعار، وضعف التنسيق اللوجستي بين المنتجين والمشترين والناقلين. يقدم هذا التقرير منصة \textbf{AgriGov Market}، وهي منصة وطنية رقمية عبر الويب صُممت من أجل رقمنة وتنظيم التجارة الزراعية المباشرة تحت الإشراف الرسمي لوزارة الفلاحة.

يربط النظام أربعة فاعلين رئيسيين — الفلاحين، المشترين، الناقلين، والوزارة — داخل منظومة رقمية موحدة. حيث يمكن للفلاحين عرض منتجاتهم وإدارتها، تنسيق عمليات النقل، وتتبع الطلبات؛ كما يمكن للمشترين تصفح سوق منظم، إجراء الطلبات، والاطلاع على سجل المشتريات؛ بينما يستطيع الناقلون قبول وتنفيذ مهام التوصيل عبر واجهة موجهة للأجهزة المحمولة؛ في حين تحتفظ الوزارة بالتحكم الإداري الكامل في التحقق من المستخدمين، إدارة تصنيفات المنتجات، وتنظيم الأسعار الرسمية.

تم تطوير النظام وفق منهجية هندسة برمجيات قياسية مقسمة إلى مرحلتين: مرحلة تحليل المتطلبات، والتي أسفرت عن مخططات حالات الاستخدام، مخططات تسلسل النظام، ونماذج أولية للواجهات؛ ومرحلة تصميم النظام، التي نتج عنها مخطط الفئات، مخطط النشاط العام، مخطط قاعدة البيانات العلائقية، وبنية النشر. تم تنفيذ المنصة باستخدام Django (واجهة REST API)، وNext.js (واجهة الويب)، وReact Native (تطبيق الهاتف)، بالإضافة إلى PostgreSQL المستضافة على Neon.

\vspace{0.5cm}
\noindent\textbf{الكلمات المفتاحية:} منصة فلاحية، سوق رقمي، تجارة مباشرة، لوجستيات، تنظيم الأسعار، Django، Next.js، React Native، وزارة الفلاحة، الجزائر.

\end{Arabic}

